\documentclass{article}
\usepackage[utf8]{inputenc}
\usepackage{graphicx}
\usepackage{amsmath}
\usepackage{geometry}
\geometry{a4paper, margin=1in}

\usepackage{accessibility}

\title{A Rust Implementation of Empirical Dependencies for Irregular Dwarf Galaxies}
\author{Jules}
\date{\today}

\begin{document}

\maketitle

\begin{abstract}
This paper presents a robust and efficient Rust implementation of the empirical formulas for irregular dwarf galaxies, as described in the work of Tillaboev, Tadjibaev, and Otojanova (2025). The implementation provides a set of functions to calculate key physical properties of these galaxies, including mass, apparent magnitude, and redshift, based on the linear regression models derived in the original study. The code is organized into a modular structure within the existing `math_explorer` framework, ensuring clarity and maintainability. Unit tests have been developed to verify the correctness of the implementation against the published formulas. This work serves as a practical computational tool for researchers in astrophysics and a demonstration of how to translate empirical scientific findings into reliable software.
\end{abstract}

\section{Introduction}

The study of dwarf galaxies is crucial for understanding the large-scale structure of the universe. These small, low-luminosity galaxies are the most numerous type of galaxy, yet they are difficult to observe. The paper "Empirical dependencies for irregular dwarf galaxies" by Tillaboev, Tadjibaev, and Otojanova (2025) makes a significant contribution to this field by establishing empirical relationships between the physical parameters of irregular dwarf galaxies.

The authors of the original study performed a statistical analysis of a catalog of 413 dwarf galaxies and derived a set of linear regression formulas that connect properties such as mass, distance, magnitude, and redshift. These formulas are valuable for predicting galaxy properties and for testing models of galaxy formation and evolution.

The goal of the work presented in this paper is to create a high-quality, well-tested Rust implementation of these empirical formulas. By encapsulating the formulas in a software library, we make them more accessible and usable for the scientific community. This implementation is part of the larger `math_explorer` project, a library of mathematical and scientific algorithms written in Rust.

\section{Methodology}

The implementation is housed within the `math_explorer` Rust crate, in a new submodule: `physics::astrophysics::galaxies`. This structure was chosen to maintain the modularity and organization of the existing project.

\subsection{Data Structures}

Two main data structures were defined to represent the concepts from the paper:

\begin{itemize}
    \item \texttt{GalaxyType}: An \texttt{enum} to represent the different morphological classifications of irregular dwarf galaxies. This allows the user to select the appropriate set of formulas for their calculations. The variants are \texttt{All}, \texttt{TypeCode10}, and \texttt{TypeCode9\_5To9\_9}.
    \item \texttt{Galaxy}: A \texttt{struct} to hold the physical properties of a galaxy. This struct is designed to be flexible, as not all properties may be known for a given galaxy. It uses Rust's \texttt{Option} type to represent optional values.
\end{itemize}

\subsection{Functions}

A set of public functions is provided to perform the calculations described in the paper. Each function corresponds to one of the empirical formulas. The function signatures are designed to be clear and self-documenting. For example:

\begin{verbatim}
pub fn calculate_log_mass_from_distance(
    distance: f64,
    galaxy_type: &GalaxyType,
) -> f64
\end{verbatim}

This function calculates the logarithm of a galaxy's mass based on its distance and morphological type. Similar functions are provided for the other relationships.

\section{Results}

To verify the correctness of the implementation, a suite of unit tests was developed. These tests compare the output of the Rust functions with the expected values calculated directly from the formulas in the paper.

For example, the paper gives the following formula for the logarithm of the mass of a galaxy of type \texttt{All} as a function of its distance:

\begin{equation}
\log(M/M_{\odot}) = 0.0230 \cdot D + 0.7840
\end{equation}

A test case was created with an input distance of 10 Mpc. The expected result is $0.0230 \cdot 10 + 0.7840 = 1.014$. The Rust function returns the same value, confirming its correctness.

The following table shows a comparison of the results from the Rust implementation and the paper's formulas for a few sample inputs.

\begin{table}[h!]
\centering
\begin{tabular}{|l|l|l|l|}
\hline
\textbf{Function} & \textbf{Inputs} & \textbf{Expected} & \textbf{Actual} \\
\hline
\texttt{log\_mass\_from\_distance} & D=10, All & 1.014 & 1.014 \\
\texttt{log\_mass\_from\_distance} & D=10, Type10 & 7.936 & 7.936 \\
\texttt{app\_mag\_from\_distance} & D=10, All & 16.207 & 16.207 \\
\texttt{log\_mass\_from\_abs\_mag} & Mv=-15, All & 8.5075 & 8.5075 \\
\texttt{redshift\_from\_log\_mass} & log(M)=10, All & -0.633 & -0.633 \\
\hline
\end{tabular}
\caption{Comparison of expected and actual results.}
\label{tab:results}
\end{table}

All tests passed, which gives us high confidence in the correctness of the implementation.

\section{Conclusion}

This paper has presented a Rust implementation of the empirical formulas for irregular dwarf galaxies from Tillaboev, Tadjibaev, and Otojanova (2025). The implementation is well-structured, documented, and thoroughly tested. It provides a valuable tool for researchers and serves as a good example of how to translate scientific research into practical software.

Future work could involve extending the library to include more models of galaxy properties, as well as incorporating the error estimates from the original paper to provide a more complete picture of the uncertainties involved in these calculations.

\begin{thebibliography}{9}

\bibitem{tillaboev2025}
K. Tillaboev, I. Tadjibaev, N. Otojanova.
\textit{Empirical dependencies for irregular dwarf galaxies}.
Astronomical and Astrophysical Transactions, 35(2), 50-62, 2025.

\end{thebibliography}

\end{document}
